\begin{figure}[t]\centering\includegraphics[width=\columnwidth]{jitter_success.pdf}\caption{Success versus start jitter on \texttt{bimodal50} (mean $\pm$ std; 5 seeds, except 10 at jitter 0.01 which are the main rows).}\label{fig:jitter}\end{figure}

\begin{table}[t]\centering\caption{Start-jitter sweep, \texttt{bimodal50}, 5 seeds (the 0.01 point is in Table~\ref{tab:main}).}\label{tab:jit}
\resizebox{\columnwidth}{!}{\begin{tabular}{lcccc}
\toprule
Method & success $\uparrow$ & collision $\downarrow$ & mode\_balance $\uparrow$ & $n$ \\
\midrule
MSE-MLP (jit=0.0) & 0.00 $\pm$ 0.00 & 1.00 $\pm$ 0.00 & 0.00 $\pm$ 0.00 & 5 \\
DDPM head (jit=0.03) & 1.00 $\pm$ 0.00 & 0.00 $\pm$ 0.00 & \textbf{0.95 $\pm$ 0.04} & 5 \\
GMM head (jit=0.03) & 1.00 $\pm$ 0.00 & 0.00 $\pm$ 0.00 & 0.94 $\pm$ 0.03 & 5 \\
MSE-MLP (jit=0.03) & 0.90 $\pm$ 0.05 & 0.10 $\pm$ 0.05 & 0.87 $\pm$ 0.06 & 5 \\
MSE-MLP (jit=0.1) & 0.98 $\pm$ 0.01 & 0.02 $\pm$ 0.01 & 0.93 $\pm$ 0.03 & 5 \\
GMM head (jit=0.1) & 1.00 $\pm$ 0.00 & 0.00 $\pm$ 0.00 & 0.94 $\pm$ 0.04 & 5 \\
DDPM head (jit=0.1) & 1.00 $\pm$ 0.00 & 0.00 $\pm$ 0.00 & \underline{0.95 $\pm$ 0.03} & 5 \\
DDPM head (jit=0.003) & 1.00 $\pm$ 0.00 & 0.00 $\pm$ 0.00 & 0.90 $\pm$ 0.05 & 5 \\
DDPM head (jit=0.0) & 1.00 $\pm$ 0.00 & 0.00 $\pm$ 0.00 & 0.95 $\pm$ 0.02 & 5 \\
GMM head (jit=0.0) & \textbf{1.00 $\pm$ 0.00} & \textbf{0.00 $\pm$ 0.00} & 0.90 $\pm$ 0.07 & 5 \\
MSE-MLP (jit=0.003) & 0.18 $\pm$ 0.05 & 0.82 $\pm$ 0.05 & 0.33 $\pm$ 0.30 & 5 \\
GMM head (jit=0.003) & \underline{1.00 $\pm$ 0.00} & \underline{0.00 $\pm$ 0.00} & 0.90 $\pm$ 0.05 & 5 \\
\bottomrule
\end{tabular}
}\end{table}

\textbf{H4: start-state jitter (supported for MSE).} MSE success is \rhval{sweep_jitter/mse-mlp-jit-0.0/bimodal50/success/mean:pct1}\% with identical starts, \rhval{sweep_jitter/mse-mlp-jit-0.003/bimodal50/success/mean:pct1}\% at 0.003, \rhval{main/mse-mlp/bimodal50/success/mean:pct1}\% at 0.01, \rhval{sweep_jitter/mse-mlp-jit-0.03/bimodal50/success/mean:pct1}\% at 0.03 and \rhval{sweep_jitter/mse-mlp-jit-0.1/bimodal50/success/mean:pct1}\% at 0.1 (Fig.~\ref{fig:jitter}, Table~\ref{tab:jit}); the difference between jitter 0.1 and 0 is \rhval{cmp/sweep_jitter/mse-mlp-jit-0.0/bimodal50/success/delta:3} ($p=\rhval{cmp/sweep_jitter/mse-mlp-jit-0.0/bimodal50/success/welch_p}$). GMM is at \rhval{sweep_jitter/gmm-head-jit-0.0/bimodal50/success/mean:pct1}\% at every value, DDPM between \rhval{sweep_jitter/ddpm-head-jit-0.0/bimodal50/success/mean:pct1}\% and \rhval{sweep_jitter/ddpm-head-jit-0.1/bimodal50/success/mean:pct1}\% (all tabulated). The MSE trend is monotone in the five tested values. A plausible reading is that the regressor can break the left/right tie using small differences in the observed position (an input-dependent mean), but we did not run an experiment that isolates this mechanism. Even at jitter 0.1 MSE stays below the mixture (MSE minus GMM difference \rhval{cmp/sweep_jitter/gmm-head-jit-0.1/bimodal50/success/delta:3}, $p=\rhval{cmp/sweep_jitter/gmm-head-jit-0.1/bimodal50/success/welch_p}$ for the difference to GMM, Welch).

\begin{table}[t]\centering\caption{DDPM denoising steps (\texttt{bimodal50}, 5 seeds). The 50-step main rows are in Table~\ref{tab:main}.}\label{tab:steps}
\resizebox{\columnwidth}{!}{\begin{tabular}{lcccc}
\toprule
Method & success $\uparrow$ & collision $\downarrow$ & mode\_balance $\uparrow$ & $n$ \\
\midrule
DDPM head (steps=1) & 0.11 $\pm$ 0.09 & 0.89 $\pm$ 0.09 & 0.04 $\pm$ 0.06 & 5 \\
DDPM head (steps=20) & \textbf{1.00 $\pm$ 0.00} & \textbf{0.00 $\pm$ 0.00} & \textbf{0.92 $\pm$ 0.06} & 5 \\
DDPM head (steps=10) & \underline{0.99 $\pm$ 0.01} & \underline{0.01 $\pm$ 0.01} & 0.90 $\pm$ 0.08 & 5 \\
DDPM head (steps=5) & 0.99 $\pm$ 0.01 & 0.01 $\pm$ 0.01 & \underline{0.91 $\pm$ 0.06} & 5 \\
\bottomrule
\end{tabular}
}\end{table}

\textbf{H5: denoising steps (partly supported).} With one step success is \rhval{abl_ddpm_steps/ddpm-head-steps-1/bimodal50/success/mean:pct1}\% (against 20 steps: difference \rhval{cmp/abl_ddpm_steps/ddpm-head-steps-1/bimodal50/success/delta:3}, $p=\rhval{cmp/abl_ddpm_steps/ddpm-head-steps-1/bimodal50/success/welch_p}$), but 5 and 10 steps are already at \rhval{abl_ddpm_steps/ddpm-head-steps-5/bimodal50/success/mean:pct1}\% and \rhval{abl_ddpm_steps/ddpm-head-steps-10/bimodal50/success/mean:pct1}\%, not distinguishable from 20 steps ($p=\rhval{cmp/abl_ddpm_steps/ddpm-head-steps-5/bimodal50/success/welch_p}$ and $\rhval{cmp/abl_ddpm_steps/ddpm-head-steps-10/bimodal50/success/welch_p}$). The hypothesis (degradation at 1 \emph{or} 5 steps) is therefore supported for 1 step only. A single ancestral step from noise may yield something close to the conditional mean, i.e.\ MSE-like behaviour; we did not test this.

\begin{table}[t]\centering\caption{GMM components (\texttt{bimodal50}, 5 seeds; $M=2$ is in Table~\ref{tab:main}).}\label{tab:comp}
\resizebox{\columnwidth}{!}{\begin{tabular}{lcccc}
\toprule
Method & success $\uparrow$ & collision $\downarrow$ & mode\_balance $\uparrow$ & $n$ \\
\midrule
GMM head (M=1) & 0.95 $\pm$ 0.02 & 0.05 $\pm$ 0.02 & 0.89 $\pm$ 0.08 & 5 \\
GMM head (M=5) & \underline{1.00 $\pm$ 0.00} & \underline{0.00 $\pm$ 0.00} & \textbf{0.96 $\pm$ 0.04} & 5 \\
GMM head (M=3) & \textbf{1.00 $\pm$ 0.00} & \textbf{0.00 $\pm$ 0.00} & \underline{0.93 $\pm$ 0.06} & 5 \\
\bottomrule
\end{tabular}
}\end{table}

\textbf{Mixture components.} A single Gaussian ($M=1$) reaches \rhval{abl_gmm_comp/gmm-head-m-1/bimodal50/success/mean:pct1}\% success (difference to $M=3$: \rhval{cmp/abl_gmm_comp/gmm-head-m-1/bimodal50/success/delta:3}, $p=\rhval{cmp/abl_gmm_comp/gmm-head-m-1/bimodal50/success/welch_p}$), well above MSE at the same jitter, probably because a wide single Gaussian sampled at each step lets the closed loop break symmetry; we did not test this. $M=2,3,5$ are all at ceiling.

\begin{table}[t]\centering\caption{Action-chunk length $K$ (open-loop execution), \texttt{bimodal50}, 5 seeds; $K=1$ is the main setting.}\label{tab:chunk}
\resizebox{\columnwidth}{!}{\begin{tabular}{lcccc}
\toprule
Method & success $\uparrow$ & collision $\downarrow$ & mode\_balance $\uparrow$ & $n$ \\
\midrule
MSE-MLP (K=4) & 0.22 $\pm$ 0.16 & 0.78 $\pm$ 0.16 & 0.43 $\pm$ 0.39 & 5 \\
MSE-MLP (K=8) & 0.03 $\pm$ 0.02 & 0.97 $\pm$ 0.02 & 0.00 $\pm$ 0.00 & 5 \\
GMM head (K=8) & \underline{1.00 $\pm$ 0.00} & \underline{0.00 $\pm$ 0.00} & \textbf{0.94 $\pm$ 0.04} & 5 \\
DDPM head (K=8) & 0.88 $\pm$ 0.01 & 0.02 $\pm$ 0.01 & 0.80 $\pm$ 0.10 & 5 \\
GMM head (K=4) & \textbf{1.00 $\pm$ 0.00} & \textbf{0.00 $\pm$ 0.00} & \underline{0.94 $\pm$ 0.06} & 5 \\
DDPM head (K=4) & 0.99 $\pm$ 0.01 & 0.00 $\pm$ 0.00 & 0.88 $\pm$ 0.10 & 5 \\
\bottomrule
\end{tabular}
}\end{table}

\textbf{Action chunk length.} Longer open-loop chunks hurt MSE badly (success \rhval{abl_chunk/mse-mlp-k-4/bimodal50/success/mean:pct1}\% at $K=4$, \rhval{abl_chunk/mse-mlp-k-8/bimodal50/success/mean:pct1}\% at $K=8$), leave GMM at \rhval{abl_chunk/gmm-head-k-8/bimodal50/success/mean:pct1}\%, and reduce DDPM to \rhval{abl_chunk/ddpm-head-k-8/bimodal50/success/mean:pct1}\% at $K=8$ (collision rate only \rhval{abl_chunk/ddpm-head-k-8/bimodal50/collision/mean:pct1}\%; the remaining rollouts time out within the 40-step horizon, whose cause we did not investigate). DDPM at $K=8$ is below DDPM at $K=4$ by \rhval{cmp/abl_chunk/ddpm-head-k-4/bimodal50/success/delta:3} and below GMM at $K=8$ by \rhval{cmp/abl_chunk/gmm-head-k-8/bimodal50/success/delta:3} ($p=\rhval{cmp/abl_chunk/gmm-head-k-8/bimodal50/success/welch_p}$). This is the one setting where the mixture beats the DDPM head here; it is a single environment with a small MLP, so it does not contradict reports that chunked diffusion works well in richer settings~\cite{chi2023diffusion}.
